\section{The Heat Death of the Human Race}

\begin{quotex}
If you love me, keep my commandments.
\flright{\textsc{John 14:15}}

Man was made in the image and likeness of God: in image he possesses freedom of will, and in likeness he possesses virtues. The likeness has been destroyed; however, man conserves the image. The image can be burnt in hell, but not consumed. It is damaged but not destroyed. Through fate as such it is not effaced, but subsists. Wherever the soul is, there also will be the image. It is not so with the likeness. This remains in the soul which accomplishes the good; in the soul which sins it is wretchedly transformed. The soul which has sinned ranks with beasts devoid of intelligence.
\flright{\textsc{St. Bernard of Clairvaux}, \emph{Sermon on the Annunciation of the Blessed Virgin Mary}}

Man is free, and remains so through all eternity --- on earth, in hell, in purgatory, in heaven --- always and everywhere. Freedom is therefore an absolute fact. As such, it entails immortality --- the argument that one finds again in Immanuel Kant's Critique of Practical Reason, for what is his ``categorical imperative'' if not the divine image in man?
\flright{\textsc{Valentin Tomberg}, \emph{Meditations on the Tarot}}
\end{quotex}


Han Fei made this observation:

\begin{quotex}
How does one address the creeping feeling that everything is getting worse and worse by the year, that the ability in people to perceive right from wrong is fading by the day, that all that is noble in man is being slowly replaced by spite, arbitrary vice and hatred.
\end{quotex}


The image to which St. Bernard refers is the reflection of the Holy Spirit in the Soul. If the Soul is quiet, that is, unperturbed by fantasies, negative emotions, and personal desires. Under those circumstances, the Soul will reflect the Spirit accurately, which is the likeness. Otherwise, the likeness is lost, much as, with fun house mirrors, the image is certainly there, but the likeness is lost.

The soul that sins is wretchedly transformed: like the souls of beasts, it is moved by concupiscence, malice, and ignorance, rather than by the pursuit of virtue. Nevertheless, as the image of God, even the sinner has freedom of will and is therefore responsible for his thoughts and acts.

\paragraph{Essential vs Accidental Evil}


\begin{quotex}
Those who enter hell are not those who have sinned accidentally, with their ``husk'' so to speak, but those who have sinned substantially or with their ``kernel'' , and this is a distinction that may not be perceptible from without; they are in any case the proud, the wicked, the hypocrites, hence all those who are the opposite of the saints and the sanctified.
\flright{\textsc{Frithjof Schuon}, \emph{Universal Eschatology}}
\end{quotex}


The existence of essential evil is beyond doubt. Essential (or substantial) evil acts are performed with full knowledge and are freely chosen. Their motivation is either the satisfaction of concupiscent desires or even deliberate maliciousness.

Accidental evil arises either from human weakness or genuine ignorance of morality. Examples of the former may be fear or peer pressure that motivates someone to do evil. Or else, the desire for some possible pleasurable experience will override the desire to do good.

In the case of ignorance, it is nevertheless a moral duty to be informed. As Tomberg points out, the moral law is written in one's soul as the divine image, so there is little excuse to remain ignorant.

Often, the evil is an unintended consequence of an act. One's liability is limited if the consequences could not be rationally anticipated. Negligence, especially reckless negligence with disregard for possible consequences, on the other hand, makes one liable for the bad consequences.

\paragraph{Essential vs accidental good}


Accidental good comes from the ``husk'' rather than from one's inner, essential nature. In the Act of Contrition this is expressed as dreading (1) the loss of heaven or (2) the pains of hell. In secular terms, one is good accidentally when the motivation is the fear of social disapproval or the threat of legal consequences.

To be essentially, or substantially good, is to follow the cosmic moral law for its own sake or for the love of God.

\paragraph{Entropy in the Social Sphere}


\begin{quotex}
In the last appearances in a cycle of decadence, of the hidden flood of spiritually destructive elementary forces after the fall of our hierarchical and traditional foundation, they are deprived of every true centre and push men and things in the most unpredictable directions.
\flright{\textsc{Julius Evola}, \emph{The Birth and Essence of the Modern Myth}}
\end{quotex}


The importance of the accidental good should not be underestimated. A healthy society will hold the ideal of the good. It will uphold this ideal through punishment if necessary. This serves to check the evil impulses of the people. The desire for social approval and the fear of social disapproval are powerful motivations.

A healthy society preserves the accumulated collective wisdom of its ancestors in its customs, habits, prejudices, legends, and myths. Otherwise, each generation would have to create itself de novo. When a society breaks down -- for a variety of reasons -- it questions that collective wisdom. While claiming to be rational, the old norms are discarded to be replaced by the latest fad. In actuality, the alleged rationality is little more than a rationalisation for the satisfaction of lower desires.

Whatever cannot be defended logically can only be defended emotionally. Lies replace truth and vice replaces virtue. The denial of the cosmic hierarchy, the moral law, leads to unpredictable consequences. The promoters of that denial are culpably and negligently liable for such consequences.

Without a hierarchy, everything becomes the same and distinctions are lost. Age old distinctions between men and women, humans and animals, saints and sinners are forgotten. That is the social equivalent of the physical principle that ``entropy is increasing''. At the end, there is only undifferentiated heat. There is certainly a lot of heat in public political discourse today.

\paragraph{The Real Good}


There is the persistent idea that the Good is equivalent to niceness. So the question always arises, ``Why does God permit evil?'' The answer is that God is pleased by the essential good and not by the accidental good. If God punished the evil and rewarded the good, then people would just be accidentally good, not essentially good. God prefers those who love Him, not those who wish to be rewarded.

Hence, the real Good is to become free, powerful, and creative. Then one is more God-like.


\flrightit{Posted on 2019-05-19 by Cologero}

\begin{center}* * *\end{center}

\begin{footnotesize}
\begin{sffamily}
\texttt{Avery on 2019-05-20 at 09:12 said:}

Anti-abortion laws will work well in conjunction with our \#MeToo culture. I think that as the spirits of these two rules synthesize, it may reverse entropy in some ways, and the children of today's pro-abortion side (if any) may well become confused as to what their parents were defending in the past.

\hfill

\texttt{Han Fei on 2019-05-20 at 15:16 said:}

So I was justified in thinking that everything ``keeps tumbling down''. Thanks I feel much better now.

But wait... this doesn't quite account the role of the elite. Evola's view, was intensely hierarchic. He saw human society, and by consequence, the worth of human beings as such, as being entirely characterized from the top down by a small group of enlightened people. My question is -- where are these enlightened people? And by where, I don't mean their physical location and identity,which is the last thing I want to know, but in terms of spirit and destiny.

When we talk about the elite, usually we mean the occluded hierarchy of those who possess esoteric knowledge and means of control, which, in my view, is ever sorely present. People are as powerless and deceived as ever by polarizing fallacies (e.g. socialism/capitalism, conservatism/liberalism, scientism/pseudoscience, atheism/false religion, etc) which the elite themselves have long put behind them. Another indication of how trivial the distinction has become between the ``elite'' and the mundane lies in how knowledge that was only disclosed to members of the upper castes and even then after a long period of mental preparation, has become openly disseminated without constraint through virulently propagating media.

Modern man's mind is bombarded by a profound multitude of symbolic images taking their forms in art, music and film, yet all of it is filtered through an impenetrable barrier of ``modernity'' installed in the mind, which renders his own grasp of how they affect his consciousness trivial and meaningless. If anything he can easily dismiss the call towards self-reflection that all these images grafted into his mind evoke by saying ``oh lol, that's totally like that trippin movie I saw with this chick that one time''. It seems that he associates all concerted effort to delve into the workings of the mind and imagination with drugs and pop culture. Who was responsible for this? In my previous post I have brought up the sinister formative influence of pop media in inverting basic notions that concern the world beyond physically felt phenomena. Yet who pushes this clearly recognizable agenda? Isn't it precisely implanted ``top down''? Society's harrowing trend towards nihilistic, self absorbed indifference can hardly be said to be of commoners' design.

So my question is, given its dismal failure in its role to serve as mankind's guide, how can we ever put any trust in any sort of hierarchy again?

\end{sffamily}
\end{footnotesize}

